%% Appendix: Proof-Pack — Independent Replay and Verification
%% wasm4pm Thesis — Periodic Table of Reason
%% Auto-generated section; treat as first-class source.

\chapter*{Proof-Pack: Independent Replay and Verification}
\addcontentsline{toc}{chapter}{Proof-Pack: Independent Replay and Verification}

\noindent
This appendix constitutes the \emph{proof-pack}: a self-contained dossier of
artefacts, pointers, and commands that allow an independent auditor to replay
every claim in the thesis from a clean checkout.  Sections A–L follow the
order of the verification pipeline: identity → registry → corpus → generation
→ refusal → OCEL replay → receipt schema → signatures → determinism →
benchmarks → replay instructions → falsifier conditions.

% -------------------------------------------------------------------------
\section*{A\quad Repository and Commit Identity}
\label{sec:repository-structure}
\addcontentsline{toc}{section}{A — Repository and Commit Identity}

\textbf{Location:} repository root (\texttt{/Users/sac/wasm4pm/}).

The canonical commit is recorded in \texttt{proof-pack/commit.txt} (generated
at pack-seal time).  To verify the exact commit against a remote:

\begin{verbatim}
git -C /Users/sac/wasm4pm log -1 --format="%H %ai %s"
git -C /Users/sac/wasm4pm verify-commit HEAD   # if GPG-signed
\end{verbatim}

\noindent\textbf{Status:} \texttt{proof-pack/commit.txt} is written by
\texttt{proof-pack/seal.sh} on every pack seal.  GPG signing is optional in
the current workflow; the SHA alone is sufficient for replay identity.

% -------------------------------------------------------------------------
\section*{B\quad Certified Breed Registry Snapshot}
\addcontentsline{toc}{section}{B — Certified Breed Registry Snapshot}

\textbf{Authoritative files:}
\begin{itemize}
  \item \texttt{breeds/registry.json} — machine-readable registry of all
        certified breeds with \texttt{id}, \texttt{category}, \texttt{tier},
        and oracle counts.
  \item \texttt{docs/registry/certified-breeds-2026-06.md} — human-readable
        snapshot as of 2026-06, listing per-breed certification status and
        open issues.
\end{itemize}

\begin{verbatim}
jq '.breeds | length' /Users/sac/wasm4pm/breeds/registry.json
wc -l /Users/sac/wasm4pm/docs/registry/certified-breeds-2026-06.md
\end{verbatim}

\noindent\textbf{Status:} Registry exists.  Snapshot document exists as of
2026-06-09.  Re-run \texttt{proof-pack/snapshot-registry.sh} to regenerate
after new breed certifications.

% -------------------------------------------------------------------------
\section*{C\quad Test Corpus Summary}
\addcontentsline{toc}{section}{C — Test Corpus Summary}

\textbf{File:} \texttt{proof-pack/oracle-counts.json}

This JSON document records, per breed, the counts of: positive oracles,
negative oracles, property-based tests generated, and OCEL replay traces.
It is produced by \texttt{proof-pack/count-oracles.sh}.

\begin{verbatim}
cat /Users/sac/wasm4pm/proof-pack/oracle-counts.json | \
  jq '{total_positive: [.[].positive] | add,
       total_negative: [.[].negative] | add}'
\end{verbatim}

\noindent\textbf{Status:} File exists; counts are updated on each CI run.
Breeds with zero oracle counts are flagged \texttt{"status": "pending"} in
the registry.

% -------------------------------------------------------------------------
\section*{D\quad Property Generator Summary}
\addcontentsline{toc}{section}{D — Property Generator Summary}

\textbf{Status:} \textsc{pending} (\(T_{\text{generated}}\) not yet wired).

The intended approach is fast-check (TypeScript) + proptest (Rust) driving the
WASM boundary.  For each breed, a property generator will:
\begin{enumerate}
  \item Sample arbitrary valid \texttt{BreedInput} via a breed-specific
        \texttt{Arbitrary} instance.
  \item Assert \texttt{status === 'ok'} and receipt non-emptiness.
  \item Assert idempotency: two runs on the same seed produce bit-identical
        \texttt{output\_hash}.
\end{enumerate}
Target: \(\geq 500\) generated cases per breed.  Generator scaffolding lives
in \texttt{packages/testing/src/generators/} (directory structure prepared;
breed impls pending).

% -------------------------------------------------------------------------
\section*{E\quad Negative / Refusal Corpus}
\addcontentsline{toc}{section}{E — Negative / Refusal Corpus}

\textbf{File:} \texttt{crates/wasm4pm-cognition/tests/oracle\_negative.rs}

Each test in this file presents a structurally malformed or semantically
illegal \texttt{BreedInput} and asserts that the WASM boundary returns
\texttt{status !== 'ok'} (or panics via \texttt{Err}).  Categories include:
missing required fields, unknown breed ids, out-of-range certainty values,
and empty premise arrays.

\begin{verbatim}
cargo test --manifest-path \
  /Users/sac/wasm4pm/crates/wasm4pm-cognition/Cargo.toml \
  oracle_negative -- --nocapture 2>&1 | grep -E "test .* ok|FAILED"
\end{verbatim}

\noindent\textbf{Status:} File exists.  Coverage is partial; breeds added
after 2026-05 require additional negative cases.

% -------------------------------------------------------------------------
\section*{F\quad OCEL Corpus and Replay Commands}
\addcontentsline{toc}{section}{F — OCEL Corpus and Replay Commands}

\textbf{Directory:} \texttt{ocel/}

\noindent\textbf{L0 traces} (exists): synthetic per-breed OCEL 2.0 JSON files
covering the 10 initially certified breeds.  Each file encodes a single lawful
execution plus a declared conforming Petri-net model.

\noindent\textbf{L1 traces} (pending): real-world XES/OCEL sourced from
\texttt{\textasciitilde/chatmangpt/pm4py} bench data — needed for the
remaining breeds.

\begin{verbatim}
# Replay all L0 traces:
bash /Users/sac/wasm4pm/proof-pack/replay-ocel.sh --level L0
# Check conformance fitness (requires pm4py):
python3 /Users/sac/wasm4pm/proof-pack/check-conformance.py ocel/L0/
\end{verbatim}

% -------------------------------------------------------------------------
\section*{G\quad Receipt Schema and Sample Receipts}
\label{app:receipt-schema}
\addcontentsline{toc}{section}{G — Receipt Schema and Sample Receipts}

\textbf{Schema:} \texttt{schemas/receipt.schema.json} (JSON Schema draft-07).

Every \texttt{wpm run} and \texttt{wpm cognition run} emits a receipt to
\texttt{.wasm4pm/receipts/latest.json}.  Required fields: \texttt{run\_id},
\texttt{breed}, \texttt{input\_hash}, \texttt{output\_hash},
\texttt{replay\_pointer}, \texttt{options\_profile}, \texttt{timestamp\_ns}.

\begin{verbatim}
# Validate a receipt against the schema:
npx ajv validate -s schemas/receipt.schema.json \
  -d .wasm4pm/receipts/latest.json
\end{verbatim}

\noindent\textbf{Status:} Schema exists.  Sample receipts in
\texttt{proof-pack/sample-receipts/} cover three breeds for illustrative
purposes.

% -------------------------------------------------------------------------
\section*{H\quad Signature Verification Transcript}
\addcontentsline{toc}{section}{H — Signature Verification Transcript}

\textbf{Status:} Receipt signing via Ed25519 (deterministic \texttt{ActorSigner})
is implemented as of v26.6.10.  Each receipt carries a \texttt{signature} field
verifiable with the actor's public key.  Production deployment requires an
institutional keypair; the private key is excluded from the open-source release.

Verifiable receipt fields:
\begin{verbatim}
{
  "run_id":          "<uuid>",
  "input_hash":      "<blake3-hex>",
  "output_hash":     "<blake3-hex>",
  "replay_pointer":  "<first-16-of-output_hash>",
  "timestamp_ns":    <u64>,
  "signature":       "<ed25519-hex>"
}
\end{verbatim}
The script \texttt{proof-pack/verify-signatures.sh} exercises the Ed25519
verification path against the actor's public key.

% -------------------------------------------------------------------------
\section*{I\quad Determinism Double-Run Transcript}
\addcontentsline{toc}{section}{I — Determinism Double-Run Transcript}

\textbf{Script:} \texttt{proof-pack/verify-determinism.sh}

\begin{verbatim}
#!/usr/bin/env bash
# Usage: bash proof-pack/verify-determinism.sh <breed> <input.json>
set -euo pipefail
RUN1=$(wpm cognition run --breed "$1" --input "$2" --no-save \
         --format json | jq -r .output_hash)
RUN2=$(wpm cognition run --breed "$1" --input "$2" --no-save \
         --format json | jq -r .output_hash)
if [ "$RUN1" = "$RUN2" ]; then
  echo "PASS determinism: $RUN1"
else
  echo "FAIL: $RUN1 != $RUN2"; exit 1
fi
\end{verbatim}

\noindent\textbf{Status:} Script exists.  CI runs it for all 13 certified
breeds on every merge to \texttt{main}.

% -------------------------------------------------------------------------
\section*{J\quad Benchmark Report}
\addcontentsline{toc}{section}{J — Benchmark Report}

\textbf{File:} \texttt{benchmarks/queueing-threshold-report.md}

This report documents the empirical queueing-threshold measurements used in
Chapter 8 (Speed Transitions).  It records, per breed and tier, the observed
\(\lambda_{\text{critical}}\) at which the breed transitions from
\texttt{realtime} to \texttt{nearline}, and from \texttt{nearline} to
\texttt{batch}.

\begin{verbatim}
# Re-run benchmarks (requires wasm-pack build first):
bash /Users/sac/wasm4pm/benchmarks/run-queueing.sh 2>&1 | \
  tee benchmarks/queueing-threshold-report.md
\end{verbatim}

\noindent\textbf{Status:} Report exists as of 2026-06-09.  Re-run after any
algorithm or WASM boundary change.

% -------------------------------------------------------------------------
\section*{K\quad Independent Replay Instructions}
\addcontentsline{toc}{section}{K — Independent Replay Instructions}

\textbf{Script:} \texttt{proof-pack/replay.sh}

\begin{verbatim}
#!/usr/bin/env bash
# Full independent replay from a clean checkout.
set -euo pipefail
cd /Users/sac/wasm4pm

# 1. Build WASM core
cd crates/wasm4pm-cognition
wasm-pack build --target nodejs --out-dir pkg -- --features wasm
cd ../..

# 2. Install TS dependencies and build
pnpm install && pnpm build

# 3. Run full test suite
pnpm test

# 4. Replay OCEL corpus
bash proof-pack/replay-ocel.sh --level L0

# 5. Verify determinism for all breeds
bash proof-pack/verify-determinism.sh --all

# 6. Print summary
cat proof-pack/oracle-counts.json | jq .
\end{verbatim}

\noindent\textbf{Status:} Script exists; steps 1–3 are confirmed working.
Steps 4–5 depend on OCEL/determinism scripts (see §F and §I).

% -------------------------------------------------------------------------
\section*{L\quad Falsifier Conditions}
\addcontentsline{toc}{section}{L — Falsifier Conditions}

\textbf{File:} \texttt{docs/falsifiers.md}

This document enumerates the precise conditions under which the central claims
of the thesis are \emph{falsified}.  Each falsifier is keyed to a theorem in
Chapter 11.  Examples:

\begin{itemize}
  \item \textbf{F-ORA}: A breed passes all positive oracles yet produces a
        different \texttt{output\_hash} on re-run with identical input
        (falsifies Theorem~O).
  \item \textbf{F-CONF}: An OCEL replay reports fitness $<0.85$ for a breed
        declared \texttt{certified} (falsifies Theorem~C).
  \item \textbf{F-NEG}: A refusal oracle input returns \texttt{status === 'ok'}
        (falsifies the refusal completeness claim).
\end{itemize}

\begin{verbatim}
cat /Users/sac/wasm4pm/docs/falsifiers.md
\end{verbatim}

\noindent\textbf{Status:} \texttt{docs/falsifiers.md} exists and is updated on
each new theorem addition.  Any CI job that triggers a falsifier condition
opens a P0 issue automatically via \texttt{.github/workflows/falsifier-gate.yml}.
